\section{Method}
We present ResearchAgent, a system that automatically proposes research ideas with LLMs.

\subsection{LLM-Powered Research Idea Generation}
\label{sec:research-idea-gen}

We begin by formally introducing the new problem of research idea generation, followed by an explanation of how LLMs are utilized to tackle it.

\paragraph{Research Idea Generation}
The goal of the research idea generation task is to formulate new and valid research ideas, to enhance the overall efficiency of the first phase of scientific discovery. While we acknowledge that the real process by which humans conduct research is varied and complex to an extent well beyond the scope of this scientific study, we attempt to model simulacra in three systematic steps that would likely be maximally beneficial to a researcher seeking assistance from an AI system. These are namely, identifying novel research ideas, proposing methods to validate these ideas, and designing experiments to measure the success of these methods in relation to the ideas.

To accomplish the aforementioned steps, we utilize the existing literature (such as academic publications) as a primary source, which provides insights about existing knowledge along with gaps and unanswered questions\footnote{We focus on the existing literature-based idea generation by following the paradigm that a \textit{new idea} is more often than not just a new combination of old elements~\cite{young2003technique}.}. Formally, let $\mathcal{L}$ be the literature, and $\vo$ be the ideas that consist of the problem $\vp$, method $\vm$, and experiment design $\vd$, as follows: $\vo = [\vp, \vm, \vd]$ where each item consists of a sequence of tokens. Then, the idea generation model $f$ can be represented as follows: $\vo = f(\mathcal{L})$, which is further decomposed into three submodular steps: $\vp = f(\mathcal{L})$ for identifying problems, $\vm = f(\vp, \mathcal{L})$ for developing methods, and $\vd = f(\vp, \vm, \mathcal{L})$ for designing experiments. We operationalize $f$ with LLMs, leveraging their capability to understand and generate academic text. 


\paragraph{Large Language Models}
Before describing the LLM in the context of our problem setup, let us first provide its general definition, which takes an input sequence of tokens $\vx$ and generates an output sequence of tokens $\vy$, as follows: $\vy = \texttt{LLM}_{\theta}(\mathcal{T}(\vx))$. Here, the model parameters $\theta$ are typically fixed after training, due to the high costs of further fine-tuning. In addition, the prompt template $\mathcal{T}$ serves as a structured format that outlines the context (including the task descriptions and instructions) to direct the model in generating the desired outputs.


\subsection{Knowledge-Augmented LLMs for~Research~Idea~Generation}
We now turn to our primary focus of automatically generating research ideas with LLMs. Recall that we aim to produce a complete idea consisting of the problem, method, and experiment design ($\vo = [\vp, \vm, \vd]$), while using the existing literature $\mathcal{L}$ as a primary source of information. We operationalize this with LLMs by instantiating the aforementioned research idea generation function $f$ with $\texttt{LLM}$ coupled with the task-specific template. Formally, $\vp = \texttt{LLM}(\mathcal{T}_p(\mathcal{L}))$ indicates the problem identification step, followed by $\vm = \texttt{LLM}(\mathcal{T}_m(\vp, \mathcal{L}))$ for method development and $\vd = \texttt{LLM}(\mathcal{T}_e(\vp, \vm, \mathcal{L}))$ for experiment design, which constitutes the full idea: $\vo = [\vp, \vm, \vd]$. 

Following this general formulation, the important question to answer is how the body of scientific literature is leveraged for actually generating research ideas with LLMs. 
Here, we outline three key desiderata that contribute to the success of human researchers ideating novel research ideas: a broad and deep understanding of related work, an encyclopedic perspective on the interconnectedness of concepts within and across scientific domains, and a community of peers who help iteratively improve ideas through constructive critiques. We describe our operationalization of these three desiderata using the prior literature and LLMs in what follows.

\paragraph{Citation Graph-based Literature Survey}
Due to the constraints on their input lengths and their reasoning abilities, particularly over very long contexts~\cite{LostITM}, it is not possible to incorporate all the existing publications from the literature $\mathcal{L}$ into the LLM input. Instead, we need to find a meaningful subset relevant to the problem at hand. To achieve this, we mirror the process followed by human researchers, who expand their knowledge of a paper by perusing other papers that either cite or are cited by it. Concretely, for the $\texttt{LLM}$, we initiate its literature review process by providing a core paper $l_0$ from $\mathcal{L}$ and then selectively incorporating subsequent papers $\left\{ l_1, ..., l_n \right\}$ that are directly connected based on a citation graph. This procedure makes the $\texttt{LLM}$ input for idea generation more manageable and coherent. In addition, we operationalize the selection process of the core paper and its relevant citations with two design choices: 1) the core paper is selected based on its citation count (e.g., exceeding 100 over 3 months) typically indicating high impact; 2) its relevant papers (which may be potentially numerous) are further narrow-downed based on their similarities of abstracts with the core paper, ensuring a more focused and relevant set of related work.

However, despite the simplicity and intuitiveness of this idea generation approach, there exists one major limitation. This approach relies exclusively on a set of given papers (the core paper and its citations); however, since scientific knowledge is not confined to specific studies but rather accumulates across a wide range of publications (across various fields), we should ideally harness this extensive, interconnected, and relevant scientific knowledge in our method for research idea generation.

\paragraph{Entity-Centric Knowledge Augmentation}
In order to model an encyclopedic view of interconnected concepts, we must effectively design a framework to extract, store and effectively leverage the vast amount of knowledge in scientific literature $\mathcal{L}$. In this work, we view entities as the atomic units of knowledge, which allows for ease of representation and accumulation over papers in a unified manner across different disciplines. For example, we can easily extract the term ``database'' whenever it appears in any paper, using existing off-the-shelf entity linking methods and then aggregate their linked occurrences into a knowledge store. Then, if the term ``database'' is prevalent within the realm of medical science but less so in hematology (which is a subdomain of medical science), the constructed knowledge store can capture the affinity between those two domains based on overlapping entities. This representational paradigm can then be used to suggest the term ``database'' when formulating the ideas about hematology. In other words, this approach enables providing novel and interdisciplinary insights by leveraging the interconnectedness of entities across various fields.

Formally, we design the knowledge store as a two-dimensional matrix $\mathcal{K} \in \mathcal{R}^{m \times m}$ where $m$ is the total number of unique entities identified and $\mathcal{K}$ is implemented in a sparse format. This knowledge store is constructed by extracting entities over all the available scientific articles in literature $\mathcal{L}$\footnote{As extracting entities on all articles is computationally infeasible, we target papers appearing after May 01, 2023.}, which not only counts the co-occurrences between entity pairs within individual papers but also quantifies the count for each entity. Our approach is versatile, thus, we can use any entity linker; in this paper we use one developed by~\citet{blink}. This off-the-shelf system proves capable of extracting key scientific entities (which is demonstrated in Table~\ref{tab:examples}) despite its lack of customized training for the scientific domain. Specifically, this linker tags and canonicalizes entities in a paper $l$ from $\mathcal{L}$, formalized as follows: $\mathcal{E}_l = \texttt{EL}(l)$ where $\mathcal{E}_l$ denotes a multiset of entities (allowing for repetitions) appearing in $l$\footnote{Due to the extensive length of scientific publications, the target of entity extraction is restricted to titles and abstracts.}. Upon extracting entities $\mathcal{E}$, to store them into the knowledge store $\mathcal{K}$, we consider all possible pairs of $\mathcal{E}$ represented as follows: $\left\{ e_i, e_j \right\}_{(i, j) \in \mathcal{C}(|\mathcal{E}|, 2)}$ where $e \in \mathcal{E}$.

Given this knowledge store $\mathcal{K}$, our next goal is to enhance the previous vanilla research idea generation process implemented based on a group of interconnected papers, denoted as follows: $\vo = \texttt{LLM}(\mathcal{T}(\left\{ l_0, l_1, ..., l_n \right\}))$. We do this by augmenting the $\texttt{LLM}$ with the relevant entities from $\mathcal{K}$, which expand the context that LLMs consume with additional knowledge. Formally, let us define entities extracted from the group of interconnected papers, as follows: $\mathcal{E}_{\left\{ l_0, ..., l_n \right\}} = \bigcup_{i=0}^{n} \texttt{EL}(l_i)$. Then, the probabilistic form of retrieving the top-$k$ relevant external entities can be represented as follows: 
\begin{equation}
\fontsize{8pt}{8pt}\selectfont
    \texttt{Ret}(\left\{ l_0, ..., l_n \right\}; \mathcal{K}) = \argmax_{I \subset [m]: |I| = k } \prod {P(e_i | \mathcal{E}_{\left\{ l_0, ..., l_n \right\}})},
    \label{eq:retrieval}
\fontsize{8pt}{8pt}\selectfont
\end{equation}
where $ [m] = \{1, ..., m\} $ and $e_i \notin \mathcal{E}_{\left\{ l_0, ..., l_n \right\}}$. Also, for simplicity, by applying Bayes' rule and assuming that entities are independent, the retrieval operation (Equation~\ref{eq:retrieval}) can be approximated as follows:
\begin{equation}
\fontsize{8.5pt}{8.5pt}\selectfont
    \argmax_{I \subset [m]: |I| = k } \prod ({\prod_{e_j \in \mathcal{E}_{\left\{ l_0, ..., l_n \right\}}} P(e_j | e_i)}) \times P(e_i),
    \label{eq:retrieval:detail}
\fontsize{8.5pt}{8.5pt}\selectfont
\end{equation}
where $P(e_j | e_i)$ and $P(e_i)$ can be derived from values in the two-dimensional matrix $\mathcal{K}$, suitably normalized. We note that the formulation in Equation~\ref{eq:retrieval:detail} is only one instance of operationalizing retrieval; this could be replaced with other retrieval strategies -- for example, embedding-based retrieval (discussions and results are provided in Appendix~\ref{appendix:retrieval}). Hereafter, the instantiation of research proposal generation augmented with relevant entity-centric knowledge is formalized as follows: $\vo = \texttt{LLM} (\mathcal{T}(\left\{ l_0, ..., l_n \right\}, \texttt{Ret}(\left\{ l_0, ..., l_n \right\}; \mathcal{K})))$\footnote{There may be additional knowledge sources (beyond the existing literature and entities) for research idea generation, and we leave exploring them as future work.}. We call this knowledge-augmented LLM-powered idea generation approach \emph{ResearchAgent}, and provide the templates to instantiate it in Tables~\ref{tab:prompt:problem}, \ref{tab:prompt:method}, and \ref{tab:prompt:experiment}.

\paragraph{Iterative Research Idea Refinements}
We note that attempting to write a full research idea in one go may not be an effective strategy. Humans write drafts that are continually improved based on multiple rounds of reviews and feedback. Therefore, we lastly model a community of peers for iterative idea improvement by introducing a set of LLM-powered reviewing agents (called \emph{ReviewingAgents}), which provide the ResearchAgent with reviews and feedback according to various criteria for improvement.

Specifically, similar to our approach to instantiate ResearchAgent with an LLM ($\texttt{LLM}$) and template ($\mathcal{T}$), ReviewingAgents are instantiated similarly but with different templates (See Tables~\ref{tab:prompt:problem:valid}, \ref{tab:prompt:method:valid}, and \ref{tab:prompt:experiment:valid}). Then, with ReviewingAgents, each of the generated research ideas (problem, method, and experiment design) is separately evaluated according to its own specific five criteria\footnote{We select the top five criteria which we consider as the most important, and leave exploring others as future work.}, which are provided in labels of Figure~\ref{fig:main} and detailed in Table~\ref{tab:criteria}. Based on the reviews and feedback from ReviewingAgents, the ResearchAgent iteratively updates and refines its generation of research ideas.

Despite the proficiency of LLMs in the evaluation of machine-generated texts~\cite{LLM/Judge/1, LLM/Judge/2}, their judgments on research ideas may not be aligned with the judgments of real human researchers. However, there are no ground truth reference judgments available, and collecting them to align LLMs is expensive and often infeasible. Ideally, the judgments made by LLMs should be similar to the ones made by humans, and we aim to ensure this by automatically generating human preference-aligned evaluation criteria (used for automatic evaluations) with a few human annotations. Specifically, to obtain these human-aligned evaluation criteria, we first collect 10 pairs of research ideas and their associated scores for every evaluation criterion on a 5-point Likert scale, annotated by human researchers having at least 3 papers. After that, we prompt the LLM with these human-annotated pairs to induce detailed descriptions for evaluation criteria~\cite{lin2024interpretable} (See Tables~\ref{tab:criteria:problem},~\ref{tab:criteria:method}, and~\ref{tab:criteria:experiment}). These criteria reflect the underlying human preferences\footnote{We additionally ask five human annotators, who evaluate research ideas, to judge the quality of the induced criteria; two of them agree strongly, while the other three agree moderately.} and are used as evaluation criteria by the ReviewingAgents.

